\documentclass{fiche}
\metadonnees{6}{Départ}{Bloc 2 — Ce qui est possible : futur et modaux}

\begin{document}
\entete

\objectifs{%
\textbf{Langue} : expression du futur ; présent après \emph{when}, \emph{if},
\emph{as soon as} ; lexique de la mer.\par
\textbf{Texte} : Robert Louis Stevenson, \emph{Treasure Island} (1883), chapitre \textsc{vii}.\par
\textbf{Durée} : 1\,h\,30 — exercices 35 min, texte et questions 30 min, version 25 min.}

%% ===================================================================
\exercice[14 min]{Dire l'avenir}

\rappel{\textbf{Rappel.} \emph{Will} annonce ou décide au moment où l'on parle ;
\emph{be going to} part d'une intention déjà prise ou d'un indice visible ;
\emph{be + -ing} dit un arrangement déjà fixé.}

\consigne{Complétez avec \emph{will}, \emph{be going to} ou le présent en \emph{be + -ing}.}

\begin{anglais}
\begin{enumerate}[label=\arabic*.,itemsep=2.2mm,leftmargin=*]
  \item ``When do we sail?'' --- ``We \trou{are sailing} tomorrow at six.'' \emph{(l'heure est
    fixée)}
  \item Look at those clouds --- there \trou{is going to be} a storm.
  \item ``Your trunk is heavy.'' --- ``Don't worry, I \trou{will carry} it for you.''
  \item The squire has bought the ship; he \trou{is going to} hire a crew next week.
  \item I \trou{am having} dinner with the doctor on Thursday; it is all arranged.
  \item ``I am afraid of the sea.'' --- ``You \trou{will get} used to it, I promise.''
  \item Take your coat: it \trou{is going to} rain before dusk.
  \item One day this boy \trou{will be} a sailor like his father.
  \item They \trou{are leaving} at dawn: the tide turns at five.
  \item I have decided: I \trou{am going to} tell my mother tonight.
  \item Nobody knows what \trou{will happen} on that island.
  \item ``Someone is at the door.'' --- ``I \trou{will go}.''
\end{enumerate}
\end{anglais}

\note{\emph{Will} : 3, 6, 8, 11, 12. \emph{Be going to} : 2, 4, 7, 10.
\emph{Be + -ing} : 1, 5, 9.}

\note[Le partage]{Décision prise sur le moment (3, 12), prédiction sans indice (8, 11),
promesse (6) : \emph{will}. Intention arrêtée avant de parler (4, 10) ou indice présent qui
annonce la suite (2, 7) : \emph{be going to}. Rendez-vous déjà pris, avec heure ou lieu (1,
5, 9) : \emph{be + -ing}. La frontière entre les trois n'est pas étanche, mais ces cas-là
sont nets.}

%% ===================================================================
\exercice[8 min]{Après \emph{when}, \emph{if}, \emph{as soon as}}

\rappel{\textbf{Rappel.} Dans une subordonnée de temps ou de condition, l'anglais n'emploie
pas \emph{will} : le présent suffit à dire l'avenir.}

\consigne{Mettez le verbe entre parenthèses à la forme qui convient.}

\begin{anglais}
\begin{enumerate}[label=\arabic*.,itemsep=2.2mm,leftmargin=*]
  \item As soon as the doctor \emph{(arrive)} \trou{arrives}, we shall sail.
  \item When I \emph{(be)} \trou{am} a man, I will go to sea.
  \item If the wind \emph{(not change)} \trou{does not change}, they will reach the island in
    a month.
  \item I will wait here until you \emph{(come)} \trou{come} back.
  \item He will send a ship after us if we \emph{(not turn up)} \trou{do not turn up} by
    August.
  \item Before the sun \emph{(rise)} \trou{rises}, the crew will be aboard.
  \item I do not know when he \emph{(come)} \trou{will come} back. \emph{(question
    indirecte)}
  \item She will put up with the noise as long as it \emph{(not last)} \trou{does not last}
    all night.
  \item When you \emph{(see)} \trou{see} the beacon, you will know we are near.
  \item Tell me if the boy \emph{(stray)} \trou{strays} too far from the quay.
\end{enumerate}
\end{anglais}

\note[Le piège de la phrase 7]{\emph{When} n'introduit pas toujours une subordonnée de temps.
Dans \emph{I do not know when he will come back}, \emph{when} introduit une question
indirecte : le futur y est normal. Test : si l'on peut remplacer \emph{when} par
\emph{at the time that}, c'est une subordonnée de temps, et le présent s'impose.}

%% ===================================================================
\exercice[13 min]{La mer et le navire}
\consigne{a) Associez chaque mot à sa traduction.}

\begin{multicols}{2}
\begin{enumerate}[label=\arabic*.,itemsep=1.2mm,leftmargin=*]
  \item a crew \hfill \trou{f}
  \item a quay \hfill \trou{i}
  \item a cove \hfill \trou{c}
  \item aboard \hfill \trou{l}
  \item a mate \hfill \trou{a}
  \item a deck \hfill \trou{g}
  \item to sail \hfill \trou{b}
  \item a sailor \hfill \trou{j}
  \item the shore \hfill \trou{d}
  \item ashore \hfill \trou{h}
  \item a tar \hfill \trou{k}
  \item a berth \hfill \trou{e}
\end{enumerate}
\columnbreak
\begin{enumerate}[label=\alph*.,itemsep=1.2mm,leftmargin=*]
  \item un second (officier)
  \item appareiller, naviguer
  \item une crique
  \item le rivage
  \item une couchette ; un poste, un emploi
  \item un équipage
  \item un pont (de navire)
  \item à terre
  \item un quai
  \item un marin
  \item le goudron
  \item à bord
\end{enumerate}
\end{multicols}

\note[Deux familles]{\emph{Aboard}, \emph{ashore}, \emph{aloft}, \emph{ahead} sont formés
sur le vieux \emph{on} : \emph{on board}, \emph{on shore}. Le préfixe \emph{a-} marque la
position. \emph{A berth} est d'abord la couchette, puis par extension la place qu'on occupe,
d'où \emph{a good berth as cook} : \og une bonne place de cuisinier \fg.}

\consigne{b) Complétez avec un mot de la liste.}
\begin{anglais}
\begin{enumerate}[label=\arabic*.,itemsep=2mm,leftmargin=*]
  \item The schooner lay at anchor, her \trou{crew} already \trou{aboard}.
  \item We walked along the \trou{quays}, among ships of every size.
  \item The old \trou{sailor} had lost his health \trou{ashore} and wanted to go to sea again.
  \item The smell of \trou{tar} and salt was something new to me.
\end{enumerate}
\end{anglais}

%% ===================================================================
\newpage
\section{Texte}

\noindent\small\textit{Jim Hawkins, quinze ans, aide sa mère à tenir une auberge au bord de
la mer. Le châtelain Trelawney vient de lui écrire qu'un navire est armé pour aller chercher
un trésor, et qu'on l'attend à Bristol.}\normalsize

\begin{texte}
You can imagine the excitement into which that letter put me. I was half beside myself with
\gl[glee]{glee}{la jubilation}; and if ever I despised a man, it was old Tom Redruth, who
could do nothing but grumble and \gl[to lament]{lament}{se lamenter}. Any of the
under-gamekeepers would gladly have changed places with him; but such was not the
\gl[a squire]{squire}{le châtelain, le hobereau}'s pleasure, and the squire's pleasure was
like law among them all. Nobody but old Redruth would have dared so much as even to grumble.

The next morning he and I set out on foot for the Admiral Benbow, and there I found my mother
in good health and spirits. The captain, who had so long been a cause of so much discomfort,
was gone where the wicked cease from troubling. The squire had had everything repaired, and
the public rooms and the sign repainted, and had added some furniture --- above all a
beautiful armchair for mother in the bar. He had found her a boy as an
\gl[an apprentice]{apprentice}{un apprenti} also so that she should not want help while I was
gone.

It was on seeing that boy that I understood, for the first time, my situation. I had thought
up to that moment of the adventures before me, not at all of the home that I was leaving; and
now, at sight of this clumsy stranger, who was to stay here in my place beside my mother, I
had my first attack of tears. I am afraid I \gl[to lead somebody a dog's life]{led that boy a
dog's life}{mener la vie dure à quelqu'un}, for as he was new to the work, I had a hundred
opportunities of setting him right and putting him down, and I was not slow to profit by
them.

The night passed, and the next day, after dinner, Redruth and I were \gl[afoot]{afoot}{à
pied, en route} again and on the road. I said good-bye to Mother and the
\gl[a cove]{cove}{une crique} where I had lived since I was born, and the dear old Admiral
Benbow --- since he was repainted, no longer quite so dear. One of my last thoughts was of
the captain, who had so often \gl[to stride, strode]{strode}{marcher à grands pas} along the
beach with his \gl[a cocked hat]{cocked hat}{un tricorne}, his sabre-cut cheek, and his old
brass telescope. Next moment we had turned the corner and my home was out of sight.

[\ldots]

While I was still in this delightful dream, we came suddenly in front of a large inn and met
Squire Trelawney, all dressed out like a sea-officer, in stout blue cloth, coming out of the
door with a smile on his face and a capital imitation of a sailor's walk.

``Here you are,'' he cried, ``and the doctor came last night from London. Bravo! The ship's
company complete!''

``Oh, sir,'' cried I, ``when do we sail?''

``Sail!'' says he. ``We sail tomorrow!''
\end{texte}
\source{Robert Louis Stevenson, \emph{Treasure Island}, Londres, Cassell, 1883,
ch.~\textsc{vii}.}

\partiel[15 min]{Questions}
\consigne{\en{Answer in English, in full sentences.}}

\begin{enumerate}[label=\arabic*.,itemsep=2mm,leftmargin=*]
  \item \en{How does Jim feel when he reads the letter, and how does Redruth feel?}
    \lignes{2}
    \corr{Jim is half beside himself with glee. Redruth does nothing but grumble and lament,
    although any of the under-gamekeepers would gladly have gone in his place.}
  \item \en{What has the squire done at the inn while Jim was away?}
    \lignes{2}
    \corr{He has had everything repaired, the public rooms and the sign repainted, added
    furniture including an armchair for Jim's mother, and found a boy to help her.}
  \item \en{Why does the sight of the new boy change everything for Jim?}
    \lignes{3}
    \corr{Until then he had thought only of the adventure ahead. The boy who will take his
    place beside his mother makes him see what he is leaving, and he cries for the first
    time.}
  \item \en{How does Jim treat the new boy, and how does he judge his own behaviour?}
    \lignes{2}
    \corr{He leads him a dog's life, correcting and humiliating him at every opportunity. The
    older narrator judges himself: \emph{I am afraid I led that boy a dog's life}.}
  \item \en{``The dear old Admiral Benbow --- since he was repainted, no longer quite so
    dear.'' Explain.}
    \lignes{2}
    \corr{The inn has been improved, but the improvement has made it a little strange to him.
    What he loves is not the building but the old familiar one, and the squire's money has
    already begun to take his home away.}
  \item \en{Find in the last lines the words that show the squire is playing a part.}
    \lignes{2}
    \corr{\emph{all dressed out like a sea-officer}, \emph{a capital imitation of a sailor's
    walk}. He is a landsman dressed up as a seaman.}
\end{enumerate}

\note[Le futur dans le texte]{\emph{When do we sail?} --- \emph{We sail tomorrow!} : présent
simple à valeur de futur, réservé à ce qui est programmé, comme un horaire
(\emph{the train leaves at six}). \emph{who was to stay here in my place} : \emph{be to} +
base verbale, futur de destin ou de programme, fréquent dans le récit.}

%% ===================================================================
\newpage
\section{Version}
\consigne{Traduisez le passage en français. Il prend place entre les deux morceaux du texte :
Jim vient d'arriver à Bristol.}

\begin{version}
Mr. Trelawney had taken up his residence at an inn far down the docks to superintend the work
upon the schooner. Thither we had now to walk, and our way, to my great delight, lay along
the quays and beside the great multitude of ships of all sizes and
\gl[a rig]{rigs}{un gréement} and nations. In one, sailors were singing at their work, in
another there were men \gl[aloft]{aloft}{en haut, dans la mâture}, high over my head, hanging
to threads that seemed no thicker than a spider's. Though I had lived by the shore all my
life, I seemed never to have been near the sea till then. The smell of tar and salt was
something new. I saw the most wonderful \gl[a figurehead]{figureheads}{une figure de proue},
that had all been far over the ocean. I saw, besides, many old sailors, with rings in their
ears, and \gl[whiskers]{whiskers}{les favoris} curled in ringlets, and tarry
\gl[a pigtail]{pigtails}{une natte, une queue de cheveux}, and their swaggering, clumsy
sea-walk; and if I had seen as many kings or archbishops I could not have been more
delighted.

And I was going to sea myself, to sea in a schooner, with a piping
\gl[a boatswain]{boatswain}{un maître d'équipage} and pig-tailed singing seamen, to sea,
bound for an unknown island, and to seek for buried treasure!
\end{version}
\source{\emph{Ibid.}}

\lignes{22}

\corr{\textbf{Proposition de traduction.} M. Trelawney avait pris ses quartiers dans une
auberge tout au bout des docks, afin de surveiller les travaux sur la goélette. C'est là
qu'il nous fallait maintenant nous rendre à pied, et notre chemin, pour mon plus grand
bonheur, longeait les quais et cette immense multitude de navires de toutes tailles, de tous
gréements et de toutes nations. Sur l'un, des matelots chantaient en travaillant ; sur un
autre, des hommes étaient perchés dans la mâture, bien au-dessus de ma tête, suspendus à des
fils qui ne semblaient pas plus épais que ceux d'une araignée. J'avais beau avoir vécu au
bord de la mer toute ma vie, il me semblait ne jamais en avoir approché avant ce jour-là.
L'odeur du goudron et du sel avait quelque chose de neuf. Je vis les plus merveilleuses
figures de proue, qui toutes étaient allées au loin par-delà l'océan. Je vis en outre
quantité de vieux marins, des anneaux aux oreilles, les favoris frisés en boucles et la
natte enduite de goudron, avec leur démarche fanfaronne et gauche de gens de mer ; et
j'aurais vu autant de rois ou d'archevêques que je n'en aurais pas été plus ravi.

Et moi-même j'allais prendre la mer, prendre la mer sur une goélette, avec un maître
d'équipage qui sifflait et des matelots nattés qui chantaient, prendre la mer à destination
d'une île inconnue, et partir en quête d'un trésor enfoui !}

\note[Choix de traduction 1]{\emph{Thither we had now to walk} : \emph{thither} = \og là,
vers ce lieu \fg{} (série \emph{hither, thither, whither}, aujourd'hui littéraire).
\emph{We had to walk} n'est pas \og nous avions marché \fg{} mais \og il nous fallait
aller \fg : \emph{have to} exprime l'obligation, sujet de la fiche 9.}

\note[Choix de traduction 2]{\emph{Though I had lived by the shore all my life, I seemed
never to have been near the sea} : l'infinitif parfait après \emph{seem} se déplie en
français par une complétive ou un infinitif composé (\og il me semblait ne jamais en avoir
approché \fg). \og J'avais beau avoir vécu \fg{} rend mieux \emph{though} que \og bien
que \fg, qui alourdirait.}

\note[Choix de traduction 3]{\emph{if I had seen as many kings\ldots I could not have been
more delighted} : irréel du passé doublé d'une comparaison. Le français retourne volontiers
la phrase : \og j'aurais vu autant de rois\ldots que je n'en aurais pas été plus ravi \fg.
Le \emph{ne} explétif est ici obligatoire.}

\note[Choix de traduction 4]{\emph{tarry pigtails} : \emph{tarry} est l'adjectif de
\emph{tar}, comme \emph{muddy} de \emph{mud}. Le français n'a pas l'adjectif et doit
étoffer : \og la natte enduite de goudron \fg.}

\note[Choix de traduction 5]{\emph{And I was going to sea myself} : \emph{to go to sea} =
\og prendre la mer, s'embarquer \fg, et non \og aller à la mer \fg. La reprise de
\emph{to sea} trois fois dans la phrase est un effet de crescendo : il faut la garder.}

%% ===================================================================
\partiel{Lexique à mémoriser}
\begin{itemize}[label=--,itemsep=0.8mm,leftmargin=*]\small
  \item \textbf{a crew} — un équipage
  \item \textbf{a quay} — un quai ; \textit{se prononce comme} key
  \item \textbf{aboard} — à bord ; a- + board\textit{, comme} ashore, aloft
  \item \textbf{to sail} — appareiller, naviguer ; \textit{nom} a sail \textit{« une voile »}
  \item \textbf{a schooner} — une goélette
  \item \textbf{the shore} — le rivage
  \item \textbf{a cove} — une crique
  \item \textbf{tar} — le goudron ; \textit{adjectif} tarry
  \item \textbf{glee} — la jubilation ; \textit{d'où} gleeful
  \item \textbf{to grumble} — ronchonner \textit{(fiche 1)}
  \item \textbf{to despise} — mépriser
  \item \textbf{to dare} — oser ; \textit{modal défectif :} he dare not
  \item \textbf{to set out} — se mettre en route
  \item \textbf{to take up one's residence} — s'installer
  \item \textbf{to superintend} — surveiller, diriger ; super- + intend\textit{, « regarder de haut »}
  \item \textbf{a thread} — un fil
  \item \textbf{to swagger} — se pavaner ; \textit{adjectif} swaggering
  \item \textbf{bound for} — à destination de
  \item \textbf{to seek} — chercher ; sought, sought
  \item \textbf{buried} — enfoui ; to bury\textit{, dont le} u \textit{se prononce} è
\end{itemize}

\end{document}
